%%%%%%%%%%%%%%%%%%%%%%%%%%%%% Define Article %%%%%%%%%%%%%%%%%%%%%%%%%%%%%%%%%%
\documentclass{article}
%%%%%%%%%%%%%%%%%%%%%%%%%%%%%%%%%%%%%%%%%%%%%%%%%%%%%%%%%%%%%%%%%%%%%%%%%%%%%%%

%%%%%%%%%%%%%%%%%%%%%%%%%%%%% Using Packages %%%%%%%%%%%%%%%%%%%%%%%%%%%%%%%%%%
\usepackage{geometry}
\usepackage{graphicx}
\usepackage{amssymb}
\usepackage{amsmath}
\usepackage{amsthm}
\usepackage{empheq}
\usepackage{mdframed}
\usepackage{booktabs}
\usepackage{lipsum}
\usepackage{graphicx}
\usepackage{color}
\usepackage[spanish]{babel}
\usepackage{psfrag}
\usepackage{pgfplots}
\usepackage{bm}


% \usepackage{biblatex}
\usepackage[backend=biber,style=vancouver,]{biblatex}
\usepackage{csquotes}
\usepackage{soul,color}
\addbibresource{biblio.bib}
%%%%%%%%%%%%%%%%%%%%%%%%%%%%%%%%%%%%%%%%%%%%%%%%%%%%%%%%%%%%%%%%%%%%%%%%%%%%%%%

% Other Settings

%%%%%%%%%%%%%%%%%%%%%%%%%% Page Setting %%%%%%%%%%%%%%%%%%%%%%%%%%%%%%%%%%%%%%%
\geometry{a4paper}

%%%%%%%%%%%%%%%%%%%%%%%%%% Define some useful colors %%%%%%%%%%%%%%%%%%%%%%%%%%
\definecolor{ocre}{RGB}{243,102,25}
\definecolor{mygray}{RGB}{243,243,244}
\definecolor{deepGreen}{RGB}{26,111,0}
\definecolor{shallowGreen}{RGB}{235,255,255}
\definecolor{deepBlue}{RGB}{61,124,222}
\definecolor{shallowBlue}{RGB}{235,249,255}
%%%%%%%%%%%%%%%%%%%%%%%%%%%%%%%%%%%%%%%%%%%%%%%%%%%%%%%%%%%%%%%%%%%%%%%%%%%%%%%

%%%%%%%%%%%%%%%%%%%%%%%%%% Define an orangebox command %%%%%%%%%%%%%%%%%%%%%%%%
\newcommand\orangebox[1]{\fcolorbox{ocre}{mygray}{\hspace{1em}#1\hspace{1em}}}
%%%%%%%%%%%%%%%%%%%%%%%%%%%%%%%%%%%%%%%%%%%%%%%%%%%%%%%%%%%%%%%%%%%%%%%%%%%%%%%

%%%%%%%%%%%%%%%%%%%%%%%%%%%% English Environments %%%%%%%%%%%%%%%%%%%%%%%%%%%%%
\newtheoremstyle{mytheoremstyle}{3pt}{3pt}{\normalfont}{0cm}{\rmfamily\bfseries}{}{1em}{{\color{black}\thmname{#1}~\thmnumber{#2}}\thmnote{\,--\,#3}}
\newtheoremstyle{myproblemstyle}{3pt}{3pt}{\normalfont}{0cm}{\rmfamily\bfseries}{}{1em}{{\color{black}\thmname{#1}~\thmnumber{#2}}\thmnote{\,--\,#3}}
\theoremstyle{mytheoremstyle}
\newmdtheoremenv[linewidth=1pt,backgroundcolor=shallowGreen,linecolor=deepGreen,leftmargin=0pt,innerleftmargin=20pt,innerrightmargin=20pt,]{theorem}{Theorem}[section]
\theoremstyle{mytheoremstyle}
\newmdtheoremenv[linewidth=1pt,backgroundcolor=shallowBlue,linecolor=deepBlue,leftmargin=0pt,innerleftmargin=20pt,innerrightmargin=20pt,]{definition}{Definition}[section]
\theoremstyle{myproblemstyle}
\newmdtheoremenv[linecolor=black,leftmargin=0pt,innerleftmargin=10pt,innerrightmargin=10pt,]{problem}{Problem}[section]
%%%%%%%%%%%%%%%%%%%%%%%%%%%%%%%%%%%%%%%%%%%%%%%%%%%%%%%%%%%%%%%%%%%%%%%%%%%%%%%

%%%%%%%%%%%%%%%%%%%%%%%%%%%%%%% Plotting Settings %%%%%%%%%%%%%%%%%%%%%%%%%%%%%
\usepgfplotslibrary{colorbrewer}
\pgfplotsset{width=8cm,compat=1.9}
%%%%%%%%%%%%%%%%%%%%%%%%%%%%%%%%%%%%%%%%%%%%%%%%%%%%%%%%%%%%%%%%%%%%%%%%%%%%%%%

%%%%%%%%%%%%%%%%%%%%%%%%%%%%%%% Title & Author %%%%%%%%%%%%%%%%%%%%%%%%%%%%%%%%
\title{Apuntes de Tesis de Licenciatura - 2022}
\author{Haoyun Qin}
%%%%%%%%%%%%%%%%%%%%%%%%%%%%%%%%%%%%%%%%%%%%%%%%%%%%%%%%%%%%%%%%%%%%%%%%%%%%%%%

\begin{document}
    \begin{titlepage}
        \centering
        {\bfseries\LARGE Universidad Peruana Cayetano Heredia \par}
        \vspace{1cm}
        {\includegraphics[width=0.3\textwidth]{img/logo_escudo.png}\par}
        \vspace{1cm}
        {\scshape\Large Facultad de Ciencias y Filosofía \par}
        \vspace{1cm}
        {\bfseries\scshape\Large Uso de modelos de aprendizaje automático en la clasificación de moléculas de colorante según su desempeño en las celdas solares orgánicas sensibilizadas por tintes naturales \par}
        \vspace{1cm}
        {\itshape\Large Tesis \par}
        \vfill
        {\Large Autor: \par}
        {\Large Bach. Dan Santivañez Gutarra\par}
        \vfill
        {\Large Asesora: \par}
        {\Large Dra. María Quintana Caceda \par}
        \vfill
        {\Large Julio 2022 \par}
    \end{titlepage}

    \tableofcontents
    \newpage

    \begin{abstract}
        Las celdas solares son dispositivos que convierten energía lumínica en 
        energía eléctrica útil a través de fenómenos electroquímicos. El desarrollo
        de materiales sostenibles y con mayor eficiencia es un reto de investigación
        y también ambiental, ya que resolverlo conlleva explorar miles de millones de 
        compuestos y hallar dicho tipos de materiales favorecerían el consumo de 
        energías limpias. Esta exploración requiere recursos humanos y materiales que
        muchos investigadores no pueden proporcionarse, es entonces que las simulaciones
        computacionales y el enfoque dirigido por datos cobran importancia. El uso de 
        experimentos \textit{in silico} han ... 
    \end{abstract}



    \newpage
    \section{Introducción}
    \subsection{Problemática}
    Muchas sociedades en está última década han estado esforzándose por dejar de depender de combustibles
fósiles y sus derivados debido a potenciales riesgos geopolíticos, económicos y ambientales. La energía solar
es una de las fuentes de energía alternativas de mayor interés por su fácil acceso y casi ilimitado suministro.
Sin embargo, el aprovechamiento de esta energía requiere de dispositivos que deben ser capáces de capturar, transformar 
y almacenar la mayor cantidad de energía lumínica en trabajo útil en rangos de tiempo variantes a causa del tiempo atmosférico y
los estadíos solares. La tecnología termosolar condensa el haz solar en un contenedor de un fluído de tal manera que este se caliente
y mueva un alternador eléctrico, mientras que la energía fotovoltaica convierte la radiación solar en energía electroquímica directamente 
por medio de fenómenos fotovoltaicos que se accionan gracias a materiales rececptores, transformadores y transportadores de carga.
Las centrales termosolares tienen el potencial de generar miles de kilos-vatios por hora (KWh) no obstante 
existen desventajas que limitan la difusión de su uso tales como su alto costo en la instalación, mantenimiento y distribución por cableado 
eléctrico, así como la gran cantidad de energía desperdiciada en los procesos de transferencia de calor. Estos problemas están casi ausentes
en los sistemas fotovoltaicos ya que se usan paneles y modulos que facilitan su uso, sin embargo aún es un reto equiparar la eficiencia de estos 
dispositivos con los termosolares. 

Los sistemas fotovoltaicos llevan en el mercado cerca de 60 años, desde su \hl{lanzamiento por Laboratorios Bell}. El corazón funcional de estos 
dispositivos son las celdas solares que son placas compuestas por materiales ordenados en forma de láminas y que tienen diferente composición y
estructura física. En las diferencias 



\begin{itemize}
    \item Cerca del \%80 de la población mundial viven de las importaciones 
     de combustibles fósiles son vulnerables a sufrir crisis energéticas
     debido a problemas geopolíticas y económicas.
     (IRENA)
    
    \item Las fuentes de energías 
\end{itemize}

    \subsection{Justificación}
    La referencia es citada en \cite{wen2020}
    \subsection{Antecedentes}
    En los antecedentes se

\subsubsection{Internacionales}

\subsubsection{Nacionales}
    \section{Marco Teórico}
    \subsection{Celdas solares sensibilizadas por tintes}
    \input{lib/MarcoTeorico/CSST.tex}
    \subsection{Enfoque dirigido por datos}
    \input{lib/MarcoTeorico/DataDriven.tex}
    \section{Metodología}
    \subsection{Análisis estadístico}

\subsubsection{Pruebas de normalidad}

La verificación de la normalidad de la distribución de un conjunto de datos se 
puede obtener a través de técnicas gráficas y por pruebas. 

En primer lugar, aquellas basadas en gráficas tienen como 



    \printbibliography
\end{document}